% GENERATED FILE. Edit canonical lessons, then run npm run generate:books.
% canonical-source-hash: fnv1a64:05e6356809970491
% canonical-lessons: LA-C328-flammeum, LA-C328-pronuba, LA-C328-nuptialis, LA-C328-matrona, LA-C328-divortium

\chapter{Nūptiae}
\label{ch:la-nuptiae}
\hlchaptermodality{328}

\begin{chapteropening}
\textbf{By the end of this chapter:} \emph{I can describe a Roman wedding and its end: the bridal veil, the matron of honour, of a wedding, the married woman and divorce.}

Complete the last lesson of chapter 328: I can describe a Roman wedding and its end: the bridal veil, the matron of honour, of a wedding, the married woman and divorce.
\end{chapteropening}

\section[flammeum, flammeī]{flammeum, flammeī — a bridal veil}

\label{lesson:LA-C328-flammeum}



\begin{quote}
Before the new one: say the Latin for a bride, then the Latin for marriage.
\end{quote}

\subsection*{You'll want to know: flammeum, flammeī}
\textbf{flammeum, flammeī} — \textquotedblleft{}a bridal veil\textquotedblright{}.

\begin{quote}
Spōnsa flammeum gerēbat.
\end{quote}

\begin{quote}
\emph{The bride wore a flame-coloured veil.}
\end{quote}

Neuter, from \emph{flammeus}, 'flame-coloured': the bride's orange-red veil.

\begin{cousinweb}
\textbf{flammeum} is the neuter of \emph{flammeus}, \textquotedblleft{}flame-coloured\textquotedblright{}, with \emph{vēlum}, \textquotedblleft{}veil\textquotedblright{}, understood (Lewis \& Short); \emph{flammeus} is built on \emph{flamma}, \textquotedblleft{}flame\textquotedblright{}.
\end{cousinweb}

\begin{grammarlens}[title={What you've built}]
One more word for telling, explaining and describing at B1.
\end{grammarlens}

\subsection*{Your turn}
\begin{itemize}
\raggedright
  \item \emph{Say it:} \emph{flammeum, flammeī}
  \item \emph{Say it:} \emph{flammeum, flammeī}, once more
  \item \emph{Say it:} a sentence of your own with \emph{flammeum, flammeī}
  \item \emph{Recall:} say the Latin for a bride, then the Latin for marriage, then say \emph{flammeum, flammeī} again
\end{itemize}

\begin{tcolorbox}[breakable,colback=teal!4,colframe=teal!35!black,title={Before you move on}]
What does flammeum, flammeī mean? (\textquotedblleft{}A bridal veil\textquotedblright{}.) Say it once more.
\end{tcolorbox}
\section[prōnuba, prōnubae]{prōnuba, prōnubae — a matron of honour}

\label{lesson:LA-C328-pronuba}



\begin{quote}
Before the new one: say the Latin for marriage, then the Latin for a bridal veil.
\end{quote}

\subsection*{You'll want to know: prōnuba, prōnubae}
\textbf{prōnuba, prōnubae} — \textquotedblleft{}a matron of honour\textquotedblright{}.

\begin{quote}
Prōnuba dextrās spōnsī et spōnsae iūnxit.
\end{quote}

\begin{quote}
\emph{The matron of honour joined the right hands of groom and bride.}
\end{quote}

Feminine. A woman married only once led the bride and joined the couple's hands.

\begin{cousinweb}
\textbf{prōnuba} is built on \emph{nūbō}, \textquotedblleft{}to marry\textquotedblright{}, with \emph{prō} (Lewis \& Short): she who helps on the marrying.
\end{cousinweb}

\begin{grammarlens}[title={What you've built}]
One more word for telling, explaining and describing at B1.
\end{grammarlens}

\subsection*{Your turn}
\begin{itemize}
\raggedright
  \item \emph{Say it:} \emph{prōnuba, prōnubae}
  \item \emph{Say it:} \emph{prōnuba, prōnubae}, once more
  \item \emph{Say it:} a sentence of your own with \emph{prōnuba, prōnubae}
  \item \emph{Recall:} say the Latin for marriage, then the Latin for a bridal veil, then say \emph{prōnuba, prōnubae} again
\end{itemize}

\begin{tcolorbox}[breakable,colback=teal!4,colframe=teal!35!black,title={Before you move on}]
What does prōnuba, prōnubae mean? (\textquotedblleft{}A matron of honour\textquotedblright{}.) Say it once more.
\end{tcolorbox}
\section[nūptiālis, nūptiāle]{nūptiālis, nūptiāle — of a wedding, bridal}

\label{lesson:LA-C328-nuptialis}



\begin{quote}
Before the new one: say the Latin for a bridal veil, then the Latin for a matron of honour.
\end{quote}

\subsection*{You'll want to know: nūptiālis, nūptiāle}
\textbf{nūptiālis, nūptiāle} — \textquotedblleft{}of a wedding, bridal\textquotedblright{}.

\begin{quote}
Post cēnam nūptiālem spōnsa ad domum novam ducta est.
\end{quote}

\begin{quote}
\emph{After the wedding feast the bride was led to her new home.}
\end{quote}

Two endings, from \emph{nūptiae} (taught). \emph{Cēna nūptiālis}: the wedding feast. English \emph{nuptial}.

\begin{cousinweb}
\textbf{nūptiālis} is built on \emph{nūptiae}, \textquotedblleft{}wedding\textquotedblright{} (EtymDB), and so on \emph{nūbō}, \textquotedblleft{}to marry\textquotedblright{}. English \emph{nuptial} comes from it.
\end{cousinweb}

\begin{grammarlens}[title={What you've built}]
One more word for telling, explaining and describing at B1.
\end{grammarlens}

\subsection*{Your turn}
\begin{itemize}
\raggedright
  \item \emph{Say it:} \emph{nūptiālis, nūptiāle}
  \item \emph{Say it:} \emph{nūptiālis, nūptiāle}, once more
  \item \emph{Say it:} a sentence of your own with \emph{nūptiālis, nūptiāle}
  \item \emph{Recall:} say the Latin for a bridal veil, then the Latin for a matron of honour, then say \emph{nūptiālis, nūptiāle} again
\end{itemize}

\begin{tcolorbox}[breakable,colback=teal!4,colframe=teal!35!black,title={Before you move on}]
What does nūptiālis, nūptiāle mean? (\textquotedblleft{}Of a wedding, bridal\textquotedblright{}.) Say it once more.
\end{tcolorbox}
\section[mātrōna, mātrōnae]{mātrōna, mātrōnae — a married woman, a lady of the house}

\label{lesson:LA-C328-matrona}



\begin{quote}
Before the new one: say the Latin for a matron of honour, then the Latin for of a wedding.
\end{quote}

\subsection*{You'll want to know: mātrōna, mātrōnae}
\textbf{mātrōna, mātrōnae} — \textquotedblleft{}a married woman, a lady of the house\textquotedblright{}.

\begin{quote}
Mātrōna servōs et domum cūrābat.
\end{quote}

\begin{quote}
\emph{The lady of the house looked after the slaves and the home.}
\end{quote}

Feminine. The \emph{stola} (taught) was the dress of a \emph{mātrōna}. English \emph{matron}.

\begin{cousinweb}
\textbf{mātrōna} is built on \emph{māter}, \textquotedblleft{}mother\textquotedblright{} (EtymDB): a married woman of standing. English \emph{matron} comes from it through French.
\end{cousinweb}

\begin{grammarlens}[title={What you've built}]
One more word for telling, explaining and describing at B1.
\end{grammarlens}

\subsection*{Your turn}
\begin{itemize}
\raggedright
  \item \emph{Say it:} \emph{mātrōna, mātrōnae}
  \item \emph{Say it:} \emph{mātrōna, mātrōnae}, once more
  \item \emph{Say it:} a sentence of your own with \emph{mātrōna, mātrōnae}
  \item \emph{Recall:} say the Latin for a matron of honour, then the Latin for of a wedding, then say \emph{mātrōna, mātrōnae} again
\end{itemize}

\begin{tcolorbox}[breakable,colback=teal!4,colframe=teal!35!black,title={Before you move on}]
What does mātrōna, mātrōnae mean? (\textquotedblleft{}A married woman, a lady of the house\textquotedblright{}.) Say it once more.
\end{tcolorbox}
\section[dīvortium, dīvortiī]{dīvortium, dīvortiī — a divorce; a separation}

\label{lesson:LA-C328-divortium}



\begin{quote}
Before the new one: say the Latin for of a wedding, then the Latin for a married woman.
\end{quote}

\subsection*{You'll want to know: dīvortium, dīvortiī}
\textbf{dīvortium, dīvortiī} — \textquotedblleft{}a divorce; a separation\textquotedblright{}.

\begin{quote}
Post decem annōs dīvortium fēcērunt.
\end{quote}

\begin{quote}
\emph{After ten years they divorced.}
\end{quote}

Neuter. \emph{Dīvortium facere}: to divorce. Also a place where roads or rivers part.

\begin{cousinweb}
\textbf{dīvortium} is built on \emph{dīvertō}, \textquotedblleft{}to turn aside, to separate\textquotedblright{} (EtymDB): a parting of ways. English \emph{divorce} comes from it through French.
\end{cousinweb}

\begin{grammarlens}[title={What you've built}]
One more word for telling, explaining and describing at B1.
\end{grammarlens}

\subsection*{Your turn}
\begin{itemize}
\raggedright
  \item \emph{Say it:} \emph{dīvortium, dīvortiī}
  \item \emph{Say it:} \emph{dīvortium, dīvortiī}, once more
  \item \emph{Say it:} a sentence of your own with \emph{dīvortium, dīvortiī}
  \item \emph{Recall:} say the Latin for of a wedding, then the Latin for a married woman, then say \emph{dīvortium, dīvortiī} again
\end{itemize}

\begin{tcolorbox}[breakable,colback=teal!4,colframe=teal!35!black,title={Before you move on}]
What does dīvortium, dīvortiī mean? (\textquotedblleft{}A divorce; a separation\textquotedblright{}.) Say it once more.
\end{tcolorbox}
